\section*{A fixed number of arbitrarily different tasks}

Fix an attempt cap $n\ge2$ and a number of designated tasks $r\ge1$. Each of these $r$ tasks may have its own success probability. The other $m-r$ tasks share one background probability. All probabilities may vary with $m$.

\textbf{Theorem.} There is an $M(n,r)$, independent of the probabilities, such that for every $m\ge M(n,r)$ and every $m\le k\le mn$,
\[
 E[Z\,{\bf1}_{\{S=k\}}]\le \Pr(S=k).
\]
Consequently $EZ\le\Pr(S\ge m)$, strictly if all probabilities are interior and at least one differs from the background. The proof does not give a numerical cutoff $M$.

Under the homogeneous law conditional on $S=k$, let $P_k$ be the uniform distribution of full strings and $Q_k=ZP_k$. Homogeneous calibration is a polynomial identity in the odds, so it holds on each stratum and $Q_k$ is normalized. Write $T=mn$, $f=T-k$, $\tau=f/m$, and $Y_i$ for the number of failures of designated task $i$. Under both laws the mean failure count of each task is $f/m$ by exchangeability. Let $t_i$ be the ratio of the failure odds of the $i$th task to those of the background. The likelihood change on the stratum is a positive constant times $t^Y=\prod_i t_i^{Y_i}$. We prove $D(t)=(E_{P_k}-E_{Q_k})t^Y\ge0$. At the all-ones vector both $D$ and its gradient vanish. Boundary probabilities follow by continuity in the original probabilities.

\textbf{Away from full success.} Take a sequence with $k/T\to a\in(1/n,1)$. The designated success counts under $P_k$ converge to $r$ independent copies of $B\sim{\rm Bin}(n,a)$. Under $Q_k$ the limits are independent copies of $W=1+{\rm Bin}(n-1-U,a)$ conditional on $U$, where $\Pr(U=u)=z^u/H(z)$ on $0\le u<n$, $H(z)=\sum_{u=0}^{n-1}z^u$, and $zH'(z)/H(z)=(1-a)/a$. The tilted mean increases from zero to $n-1$ because its derivative with respect to $\log z$ is the positive variance.

Indeed $Q_k$ assigns total wait $s$ mass
\[
 \binom{m+s-1}{s}\binom{T-m-s}{k-m}\big/\binom{T}{k}.
\]
This is the law of the position of the $m$th success minus $m$ in a uniform string with $k$ successes. Counting successes in prefixes of lengths $\lfloor m(1/a\pm\epsilon)\rfloor$ gives $s/m\to(1-a)/a$ in probability, since the means bracket $m$ by order $m$ and variances are $O(m)$. Conditional on $s$, waits form a uniform capped composition and suffix successes are uniformly placed. Tilt the capped sites by $z_s$ to have mean $s/m$. The probability of a prescribed pattern of $r$ waits divided by its independent tilted probability is a ratio of local sum probabilities on $m-r$ and $m$ sites. Along every sequence $s/m\to(1-a)/a$ both probabilities are asymptotic to $1/\sqrt{2\pi m\sigma_s^2}$, with $\sigma_s^2$ the tilted variance. The displacements from the means are bounded. In Fourier inversion the centered characteristic function is $1-\sigma_s^2 v^2/2+O(|v|^3)$ uniformly near zero; off a neighborhood of zero in $[-\pi,\pi]$ its modulus stays below one uniformly because the probabilities of zero and one stay positive. Gaussian scaling gives the local estimate. Sampling finitely many suffix sites gives the independent Bernoulli limits since $(k-m)/(T-m-s)\to a$. At $a=1/n$ the full limits are still binomial and each weighted success count tends to one since it is at least one and its mean tends to one.

At an interior $a$ the full one-site limit strictly dominates the weighted one for every hinge $(j-h)_+$, $1\le h<n$. Let $U_B$ be the first-hit wait in the binomial string conditional on at least one hit. For $j\ge1$ the ratio of weighted to full mass is
\[
 \frac{1}{aH(z)}E[(z/(1-a))^{U_B}\mid B=j].
\]
This follows from the common suffix law and the binomial first-hit mass $a(1-a)^u$. Truncating the geometric law lowers its mean, so $z>1-a$. Given $j$ hits, their sites form a uniform subset. The minimum position decreases with positive probability at each increment of $j$ under a random-permutation coupling. The mass ratio therefore decreases strictly in $j$. Both laws have mean $na$. Weight both by $j/(na)$; on the positive support the ratio of full to weighted mass is now strictly increasing. Its covariance with the nonconstant increasing function $(j-h)_+/j$ under the size-weighted stopped law is positive, proving strict hinge dominance. At $a=1/n$ the binomial law strictly dominates the constant one on all the same hinges. Reflection gives these inequalities for failures as well. In particular the full variance exceeds the weighted variance.

For uniformity in $t$, put $b_i=t_i/(1+t_i)$ and multiply $D(t)$ by $\prod_i(1-b_i)^n$. The resulting polynomial extends to the closed cube. Its coefficients converge along the preceding sequences. Its limit is the difference of two products of one-site expectations with factors $b_i^y(1-b_i)^{n-y}$. Each full expectation is positive. If $b_i\ne1/2$, strict hinge dominance gives strict inequality for this factor: in the open interval the factor is strictly convex in $y$, and at zero and one it is an endpoint indicator with positive hinge gap. Hence the limiting product gap is positive unless every $b_i=1/2$.

Near that one point use $t$ again. Finite support gives convergence of second derivatives on compact neighborhoods. The limiting Hessian of $D$ at the all-ones vector is a positive multiple of the identity. Diagonal entries are second factorial moment differences, equal to the positive variance gap; mixed entries vanish by independence and equality of means. In a neighborhood, Taylor's formula and the exact zero value and gradient give a positive lower bound proportional to $\|t-{\bf1}\|^2$ for sufficiently large members of the sequence. This covers changes converging arbitrarily quickly to equality.

\textbf{Near full success.} For $f<n$ the score is identically one. Take $n\le f\le\delta m$ for $\delta>0$ sufficiently small, depending only on $n,r$. A specified pattern on the $nr$ marked slots with failure vector $y$ and $d=\sum_i y_i\le f$ has full mass $(f)_d(T-f)_{nr-d}/(T)_{nr}$, between $c\tau^d$ and $C\tau^d$. Here parentheses denote falling factorials. The bounds use fixed $nr$ and $(f)_d/f^d\ge d!/d^d$ for $f\ge d\ge1$. Any pattern exhausting a task has zero weighted mass. For other patterns the ratio of weighted to full mass is $1+O_{n,r}(\tau^{n-1})$. At $d=0$ the ratio is at most one.

To prove this last comparison, let $w_j(s)$ be the reciprocal of cap compliance for a uniform composition on $j$ sites. For $m-r\le j\le m$, $s\le f$, a specified pattern on one or two sites of total $e\ge1$ under the capped law has probability at most $(Cs/j)^e$. Clear those sites and redistribute their units to distinct empty outside sites. Each source has at least $\binom{j-2-s}{e}$ images and each image at most $\binom{s}{e}$ inverse choices. Inserting a star at a uniformly chosen position of a uniform stars-and-bars word gives a uniform word. Conditional on compliance, violation occurs when the chosen part is saturated, with probability at most $C\tau^{n-1}$ by the one-site estimate (selection weight is size plus one). Deleting a uniformly chosen bar gives a uniform word with one fewer part. A compliant word fails only if the adjacent parts sum to at least $n$, with probability at most $C\tau^n$ by exchangeability and the two-site estimate. Noncompliance survives either operation. Thus
\[
 1\le w_j(s+1)/w_j(s)\le1+C\tau^{n-1},\qquad
 1-C\tau^n\le w_{j+1}(s)/w_j(s)\le1 .
\]
For a nonexhausted marked record, the remaining homogeneous stratum has at least $n(m-r)-f\ge m-r$ successes, so its own score is normalized. On its scored paths compare $w_m(s+u)$ with $w_{m-r}(s)$, where the record adds $u\le r(n-1)$ waits and $s+u\le f$. Applying the two estimates proves the mass ratio. At $d=0$ there are no added waits.

First bound $0\le t_i\le K$ for fixed $K>1$. In the signed sum for $D$ use $g_y=t^y-1-\sum_i y_i(t_i-1)$, by equality of masses and means. It vanishes at $d=0,1$. Its absolute value is at most $C_{K,n,r}\|t-{\bf1}\|^2$ by Taylor's formula. The total contribution of nonexhausted patterns is therefore bounded in absolute value by $C\tau^{n+1}\|t-{\bf1}\|^2$. The same bound holds for exhausted patterns of degree above $n$. At degree $n$ the exhausted patterns have all failures on one task. Each has full mass $\alpha(1+O(\tau))$, $\alpha=(f)_n/(T)_n$, by a union bound for additional marked failures conditional on those $n$ failures. Consequently
\[
 D(t)=\alpha\sum_{i=1}^r[t_i^n-1-n(t_i-1)]
        +O_{K,n,r}(\tau^{n+1}\|t-{\bf1}\|^2).
\]
Here $\alpha\ge c\tau^n$. Each bracket is bounded below by a positive multiple of $(t_i-1)^2$, since their quotient extends positively at one and is positive elsewhere on $[0,K]$. Taking $\delta$ small gives the sign uniformly on this box.

For $t_*=\max_i t_i\ge K$, put $v_i=\tau t_i$ and $v=\max_i v_i$. Up to fixed positive factors, the full tilted mass of each feasible labelled pattern equals $\prod_i v_i^{y_i}$. Use the original nonnegative monomials in this comparison. The zero-failure pattern contributes nonnegatively. If $v\le1$, the sum of tilted full masses of nonzero nonexhausted patterns is at most $Cv$ since their number is fixed. Their negative contribution is at worst $C\tau^{n-1}v$. Exhausting only a task with maximal odds contributes at least $cv^n$. Choosing $K$ large enough makes the latter larger since $v^{n-1}/\tau^{n-1}=t_*^{n-1}$.

If $v\ge1$, each nonexhausted pattern is dominated in monomial weight by an exhausted one. At $d\ge n$ move failures from other marked slots into the maximal task until it is exhausted, preserving $d$. At $d<n$, concentrate all failures there and add $n-d$ failures, feasible since $f\ge n$. These operations weakly increase the weight and the remaining bank has enough slots. There are only $2^{nr}$ patterns. Thus the total nonexhausted tilted mass is at most a constant times the exhausted mass, and only a factor $O(\tau^{n-1})$ of it can contribute negatively. Choose $K$ first and then $\delta$ to obtain all three bounds simultaneously.

A sequence violating the theorem outside this endpoint region would have a subsequence with $k/T$ converging in $[1/n,1-\delta/n]$ and convergent compactified odds. The product comparison and its diagonal Hessian exclude every such limit. There are interior strata for large $m$, giving strictness of the sum for unequal interior probabilities. A bank with an unbounded number of changed tasks is still outside the result.
